\citet{toyota} expands on the work of \citet{merrill2023parallelism}.
\citet{merrill2023parallelism} have shown that one step of a constant-depth, polynomial-embedding-dimension and log-precision deterministic transformer can be simulated in a small parallel time, i.e., using $\TC^0$ circuits.
This result is built on the fact that multiplication and division of $n$-bit binary numbers are in $\TC^0$ \citep{hesse2001division}, as well as the iterated addition over $n$ different $n$-bit binary integers. \bigcifu{Ok, se usa TC.}

However, such $\TC^0$ expressiveness upper bounds may be unrealistic for transformers operating with floating-point numbers. 
\citet{merrill2023parallelism} implicitly assumed that, when adding more than one floating-point number, the algorithm first computes the exact answer without rounding, using arbitrarily higher precision, and performs rounding only at the end.
However, in practice, rounding occurs after each addition of two numbers, and it remains open whether such $\TC^0$ upper bounds still hold.
Immediate rounding makes iterated addition over floating-point numbers no longer associative \citep{goldberg1991every}.
The associativity of integer addition is essential for the result that iterated addition over $n$ different $n$-bit binary integers is in $\TC^0$.

The abstraction for transformers presented by \citet{toyota} uses the numeric representation presented in Section \ref{repr}.
Their model rounds after every numerical operation and represents numbers with a constant number of bits for the significand and the exponent.

\citet{toyota} prove that $\mathsf{sum}_{e,s}:\Floating^n \to \Floating$ has an $\AC^0$ circuit when $e,s$ are both constants.
Then, composing circuits, Theorem \ref{theo:ac0-sim-det-transf} is obtainded.

\bigcifu{Sería mas extenso acá. "Notar que todas las componentes del algoritmo X se pueden implementar en AC0 dada que solo involucran operaciones sobre $F_{e,s}$, por lo que usando esta observacion se tiene que...}

\begin{theorem}\label{theo:ac0-sim-det-transf}
    One step of a deterministic, constant-depth, polynomial-embedding-dimension transformer can be simulated with a circuit in $\AC^0$.
\end{theorem}

Observe that the fact that a single forward pass of a transformer can be simulated by an $\AC^0$ circuit immediately \bigcifu{No me gusta el immediately xD, pero opinable}implies that a transformer with $O(\log(n))$ steps of $\CoT$ can also be simulated by $\AC^0$.
Thus, we obtain the following theorem.

\begin{theorem}\label{theo:det-upper-bound-log-steps}
    $\CoT[\log(n), \poly(n)] \subseteq \AC^0$
\end{theorem}

\begin{proof}
    Theorem \ref{theo:det-upper-bound-log-steps} is sustained on the fact that a transformer with $T$ steps of $\CoT$ can be interpreted as an $\mathsf{OR}$ of $|\Sigma|^T$ disjoint subcircuits, where each of them enumerates all possible sequences of $T$ chain-of-thought tokens and outputs the value of the token in the branch where all the intermediate token values are consistent. 
    
    We say a sequence of tokens $\sigma_1, \dots, \sigma_T$ is consistent with the input $\alpha$ when it satisfies $\TF(\alpha) = \sigma_1$, $\TF(\alpha \sigma_1) = \sigma_2$, $\dots$, $\TF(\alpha \sigma_1 \dots \sigma_{T-1}) = \sigma_{T}$.
    Each of the conditions can be checked by an $\AC^0$ circuit in parallel.
    Then, for each sequence of tokens there is an $\AC^0$ circuit validating its consistency.
    Since the transformer is deterministic, there is only one consistent sequence.
    The circuit outputs the value of the last token in the consistent sequence.
    
    The enumeration of all possible sequences can be done in parallel and thus only takes constant depth.
    When $T = O(\log(n))$, this only leads to a $\poly(n)$ blow-up in circuit size.
    Therefore, it follows that $\CoT[\log(n), \poly(n)] \subseteq \AC^0$.
\end{proof}

\bigcifu{Quedo muy bien esta proof.}

\medskip

Since each pass of the transformer can be simulated by a circuit in $\AC^0$, a transformer performing a polynomial number of $\CoT$ steps $T(n)$ can be simulated by a circuit of depth $T(n)$ (and polynomial size), by stacking the circuit for one pass.
Therefore, a transformer performing a polynomial number of steps can be simulated by a circuit in $\ppoly$.

\begin{corollary}\label{det-poly-steps-in-ppoly}
    $\CoT[\poly(n), \poly(n)] \subseteq \ppoly$
\end{corollary}